% GENERATED FILE. Edit canonical lessons, then run npm run generate:books.
% canonical-source-hash: fnv1a64:ee76b646980d1c27
% canonical-lessons: KA-C16-tingalugalu, KA-C16-uttara, KA-S126-letter-u, KA-S147-digit-eight, KA-C16-checkpoint

\chapter{The Traditional Months}
\label{ch:ka-traditional-months}
\hlchaptermodality{16}

\begin{chapteropening}
\textbf{By the end of this chapter:} \emph{I can name the twelve traditional Kannada months and say when the new year Ugadi falls.}

Complete KA-C16-checkpoint, the chapter's closing checkpoint: say the twelve months, place Ugadi on Chaitra Śukla Pratipada and say the calendar is shared with Hindi, recall uttara, find \kn{ಉ} in \kn{ಉತ್ತರ} and read \kn{೮} as eṇṭu, then write both from sound.
\end{chapteropening}

\section[\texorpdfstring{\kn{ಚೈತ್ರ} \kn{ವೈಶಾಖ} \kn{ಜ್ಯೇಷ್ಠ} … (caitra vaiśākha …)}{ಚೈತ್ರ ವೈಶಾಖ ಜ್ಯೇಷ್ಠ … (caitra vaiśākha …)}]{\kn{ಚೈತ್ರ} to \kn{ಫಾಲ್ಗುಣ} — the calendar Kannada shares with Hindi, not with Tamil}

\label{lesson:KA-C16-tingalugalu}



\begin{quote}
Tamil and Malayalam each built their own solar calendar. Kannada takes a different path entirely — the same path as Hindi.
\end{quote}

\subsection*{You'll want to know: The twelve months}
\textbf{\kn{ಚೈತ್ರ}, \kn{ವೈಶಾಖ}, \kn{ಜ್ಯೇಷ್ಠ}, \kn{ಆಷಾಢ}, \kn{ಶ್ರಾವಣ}, \kn{ಭಾದ್ರಪದ}, \kn{ಆಶ್ವಯುಜ}, \kn{ಕಾರ್ತೀಕ}, \kn{ಮಾರ್ಗಶಿರ}, \kn{ಪುಷ್ಯ}, \kn{ಮಾಘ}, \kn{ಫಾಲ್ಗುಣ}} (\emph{Caitra, Vaiśākha, Jyēṣṭha, Āṣāḍha, Śrāvaṇa, Bhādrapada, Āśvayuja, Kārtīka, Mārgaśira, Puṣya, Māgha, Phālguṇa}).

\begin{culture}
This is the \textbf{same pan-Indian Sanskritic lunisolar calendar} you already met as Hindi's \textbf{Vikram Samvat} — the same twelve month names, tied to the same \textbf{six-ṛtu} seasonal cycle. Karnataka's own New Year, \textbf{Ugadi}, falls on \emph{Chaitra Śukla Pratipada} — the first day of the bright fortnight of Chaitra — exactly the reasoning behind Hindu New Year celebrations across much of India. Unlike Tamil's or Malayalam's genuinely home-grown solar calendars, Kannada didn't build its own separate system here — it shares this one with Hindi, Telugu, and much of the rest of the subcontinent.
\end{culture}

\subsection*{Your turn}
\begin{itemize}
\raggedright
  \item \emph{Say it:} the twelve — \textquotedblleft{}Caitra, Vaiśākha, Jyēṣṭha, Āṣāḍha, Śrāvaṇa, Bhādrapada, Āśvayuja, Kārtīka, Mārgaśira, Puṣya, Māgha, Phālguṇa\textquotedblright{}
  \item \emph{Say it:} \textquotedblleft{}Ugadi — Chaitra Śukla Pratipada\textquotedblright{}
  \item \emph{Say it:} the honest point — shared with Hindi, unlike Tamil/Malayalam's own calendars
  \item \emph{Recall:} say \emph{tale kai}
\end{itemize}

\begin{tcolorbox}[breakable,colback=teal!4,colframe=teal!35!black,title={Before you move on}]
Does Kannada use its own solar calendar, like Tamil and Malayalam? (\textbf{No} — it uses the \textbf{same pan-Indian lunisolar calendar} as Hindi's Vikram Samvat.) What is Kannada's New Year called, and when does it fall? (\textbf{Ugadi}, on \emph{Chaitra Śukla Pratipada}.)
\end{tcolorbox}
\section[\texorpdfstring{\kn{ಉತ್ತರ} (uttara)}{ಉತ್ತರ (uttara)}]{\kn{ಉತ್ತರ} (uttara) — an answer; north}

\label{lesson:KA-C16-uttara}



\begin{quote}
Before the new word: say the last thing you learned, once, out loud.
\end{quote}

\subsection*{You'll want to know: \kn{ಉತ್ತರ}}
\textbf{\kn{ಉತ್ತರ}} — \emph{uttara} — \textquotedblleft{}an answer; north\textquotedblright{}.

Say it: \emph{uttara}. Look at its shape as you say it: the next lesson takes one piece of it, \textbf{\kn{ಉ}}, and writes it on its own.

\subsection*{Your turn}
Say these aloud:

\begin{itemize}
\raggedright
  \item \emph{uttara}
  \item \emph{uttara}, and what it means
\end{itemize}

\begin{tcolorbox}[breakable,colback=teal!4,colframe=teal!35!black,title={Before you move on}]
What is \textquotedblleft{}an answer; north\textquotedblright{}? (\textbf{\kn{ಉತ್ತರ}}, \emph{uttara}.)
\end{tcolorbox}
\section[\texorpdfstring{\kn{ಉ} (u)}{ಉ (u)}]{\kn{ಉ} — one character, the standing half of a sign you already read}

\label{lesson:KA-S126-letter-u}



\begin{quote}
Before the new one: \kn{ಪ} — what sound does it carry?

One character this time, and you are already reading half of it.
\end{quote}

\subsection*{Script you'll notice: \kn{ಉ}}
\textbf{\kn{ಉ}} — \emph{u}.

It is an \textbf{independent vowel}: the shape a vowel takes when it stands at the start of a word with no consonant in front of it.

You have been reading its other half since the chapter that taught \textbf{\kn{ು}}, the vowel \emph{sign} — the same sound, drawn as a tail that rides on the consonant before it. Look at the ending of each of these, which you can already say:

\begin{itemize}
\raggedright
  \item \textbf{\kn{ಹೌದು}} \emph{haudu} — yes
  \item \textbf{\kn{ಇರು}} \emph{iru} — to be, to stay
  \item \textbf{\kn{ನೀವು}} \emph{nīvu} — you
\end{itemize}

Every one of them ends in that riding tail. Take the consonant away and the same vowel has to stand on its own feet, and \textbf{\kn{ಉ}} is what it looks like standing.

So one sound, two shapes, and where you find it in the word tells you which one you are looking at:

\begin{itemize}
\raggedright
  \item \textbf{\kn{ಉ}} — standing alone, at the beginning.
  \item \textbf{\kn{ು}} — attached, riding on the consonant before it.
\end{itemize}

This is the same arrangement as \textbf{\kn{ಓ}} and \textbf{\kn{ೋ}}, and as \textbf{\kn{ಇ}} and \textbf{\kn{ಿ}}. Kannada does it for every vowel it has.

\subsection*{Writing: \kn{ಉ} — follow the numbered strip}
\hlblockfigure{\detokenize{figures/KA-S126-letter-u-filmstrip.pdf}}{How \kn{ಉ} is written, stroke by stroke}

Follow the numbered strip: start where movement 1 begins, and let each panel show you the next movement before your pen makes it. Then write \kn{ಉ} yourself — slowly, and larger than it is printed.

\begin{quote}
The strip shows \textbf{where to start the character and which way to travel}, in one attested order. That is taught with real variation from school to school, so the strip names its source beneath it: treat it as a sound way in, not the only one.
\end{quote}

\subsection*{Your turn}
\begin{itemize}
\raggedright
  \item \emph{Look:} at these two and say which one carries the riding tail, not the standing letter
\end{itemize}

\begin{quote}
\kn{ಹೌದು}  ·  \kn{ಉ}
\end{quote}

\begin{itemize}
\raggedright
  \item \emph{Trace it:} \kn{ಉ} three times, saying \emph{u} as you finish each one
  \item \emph{Look:} at \kn{ಉ} and \kn{ು} side by side, and say which one stands alone
\end{itemize}

\begin{tcolorbox}[breakable,colback=teal!4,colframe=teal!35!black,title={Before you move on}]
Which character is this — \kn{ಉ}? What sound does it carry? (\emph{\textbf{u}}.) What is the difference between \kn{ಉ} and \kn{ು}? (One stands alone at the start of a word; the other rides on the consonant before it.) Name a word you already say that ends in the riding one.
\end{tcolorbox}
\section[\texorpdfstring{\kn{೮} (8)}{೮ (8)}]{\kn{೮} — the digit for eṇṭu}

\label{lesson:KA-S147-digit-eight}



\begin{quote}
Before the new shape: \textbf{\kn{೭}} — which number is it?

One digit this time, and you already know its name.
\end{quote}

\subsection*{Script you'll notice: \kn{೮}}
\textbf{\kn{೮}} — \textbf{8}, the digit for \textbf{\kn{ಎಂಟು}} (\emph{eṇṭu}), \textquotedblleft{}eight\textquotedblright{}.

Eight. The two numbers whose words wander furthest from Tamil and Telugu have the least surprising digits.

\subsection*{Writing: \kn{೮} — follow the numbered strip}
\hlblockfigure{\detokenize{figures/KA-S147-digit-eight-filmstrip.pdf}}{How \kn{೮} is written, stroke by stroke}

Follow the numbered strip: start where movement 1 begins, and let each panel show you the next movement before your pen makes it. Then write \kn{೮} yourself — slowly, and larger than it is printed.

\begin{quote}
The strip shows \textbf{where to start the character and which way to travel}, in one attested order. That is taught with real variation from school to school, so the strip names its source beneath it: treat it as a sound way in, not the only one.
\end{quote}

\subsection*{Your turn}
\begin{itemize}
\raggedright
  \item \emph{Look:} at the digits you have met, and name each one aloud
\end{itemize}

\begin{quote}
\kn{೧}  ·  \kn{೨}  ·  \kn{೩}  ·  \kn{೪}  ·  \kn{೫}  ·  \kn{೬}  ·  \kn{೭}  ·  \kn{೮}
\end{quote}

\begin{itemize}
\raggedright
  \item \emph{Say it:} \textquotedblleft{}eṇṭu\textquotedblright{} as you point at \kn{೮}
\end{itemize}

\begin{tcolorbox}[breakable,colback=teal!4,colframe=teal!35!black,title={Before you move on}]
Which number is \textbf{\kn{೮}}? (\textbf{8} — \emph{eṇṭu}, \kn{ಎಂಟು}.) Name every digit you have met so far, in order.
\end{tcolorbox}
\section[Practice]{Chapter 16 checkpoint — the months, Ugadi, \kn{ಉ} and \kn{೮}}

\label{lesson:KA-C16-checkpoint}



\begin{quote}
Nothing new here. The months and the new year out loud first, then one word, one letter and one digit on the page.
\end{quote}

\subsection*{Your turn — the months}
\begin{itemize}
\raggedright
  \item \emph{Say it:} the twelve months, from the first
\end{itemize}

(\emph{Caitra, Vaiśākha, Jyēṣṭha, Āṣāḍha, Śrāvaṇa, Bhādrapada, Āśvayuja, Kārtīka, Mārgaśira, Puṣya, Māgha, Phālguṇa}.)

\begin{itemize}
\raggedright
  \item Is this calendar Kannada's own, the way Tamil's and Malayalam's are theirs? (No: it is the pan-Indian lunisolar calendar Hindi uses too.)
  \item When does Ugadi fall? (On \emph{Chaitra Śukla Pratipada}, the first day of the bright fortnight of Chaitra.)
\end{itemize}

\subsection*{Your turn — one more word from this chapter}
\begin{itemize}
\raggedright
  \item What is \textquotedblleft{}an answer\textquotedblright{}, and also \textquotedblleft{}north\textquotedblright{}? (\emph{Uttara}.)
\end{itemize}

\subsection*{Script — \kn{ಉ} and \kn{೮}}
\begin{itemize}
\raggedright
  \item \textbf{\kn{ಉತ್ತರ}} is \emph{uttara}. Which character opens it, and why does it stand on its own there? (\textbf{\kn{ಉ}}, \emph{u}: the word starts on a vowel, so the vowel takes its standing shape.)
  \item Where would the same \emph{u} take its riding shape, \textbf{\kn{ು}}? (After a consonant, as at the end of \textbf{\kn{ಹೌದು}}.)
  \item Which number is \textbf{\kn{೮}}, and which word does it stand for? (\textbf{8}, \emph{eṇṭu}.)
\end{itemize}

\emph{Read it:} \textbf{\kn{ಉತ್ತರ}} aloud, and point to the \kn{ಉ} in it

\emph{Read it:} \textbf{\kn{೮}} aloud as \emph{eṇṭu}

\subsection*{Writing — from sound}
\emph{Cover:} the Script block above, so no model stays in view

\emph{Write it:} \emph{u} as it stands at the start of a word, then the digit for \emph{eṇṭu}

\emph{Check:} against the key — \textbf{\kn{ಉ}}, \textbf{\kn{೮}}

\emph{Write it:} one repaired shape, then stop

\begin{tcolorbox}[breakable,colback=teal!4,colframe=teal!35!black,title={Before you move on}]
Which month opens the year, and which festival opens it? (\emph{Caitra}; Ugadi.) That is the chapter: a calendar shared with Hindi, a word that means both answer and north, and the vowel and the digit that sit inside it.

Next chapter: noon and midnight.
\end{tcolorbox}
